El TP x86 tenía \texttt{taskA.c}/\texttt{taskB.c} + \texttt{idle.asm}, compilados como blobs que el kernel copiaba desde
\texttt{0x10000} al tablero. Este port sigue el mismo espíritu, pero usa un \emph{build separado} para los binarios de
usuario:

\begin{itemize}
  \item Fuentes: \texttt{riscv/user/miner.c}, \texttt{riscv/user/idle.c}, \texttt{riscv/user/entry.S}.
  \item Linker script propio (\texttt{riscv/user/linker.ld}) con \texttt{.}~$=$~\texttt{0x08000000}.
  \item Producen \texttt{build/user\_miner.bin} / \texttt{build/user\_idle.bin} con \texttt{objcopy -O binary}.
  \item Se embeben en el kernel como símbolos \texttt{\_binary\_build\_user\_miner\_bin\_start} / \texttt{\_end} vía
  \texttt{ld -r -b binary}~\cite{gnu-ld}. El kernel los copia a la celda correspondiente en \texttt{write\_cell\_image}
  durante \texttt{mmu\_task\_init}.
\end{itemize}

El \texttt{entry.S} de usuario es trivial:

\begin{codesnippet}
\begin{verbatim}
 _start:
     .option push
     .option norelax
     la gp, __global_pointer$
     .option pop
     call user_main
 1:  j 1b
\end{verbatim}
\end{codesnippet}

El ABI para invocar servicios del sistema está en \texttt{riscv/syscall.h}, usando la convención \texttt{a7}=id +
\texttt{a0}=arg/retorno (estándar \emph{syscall} de Linux/RISC-V~\cite{riscv-psabi}):

\begin{codesnippet}
\begin{verbatim}
 static inline void syscall_move(e_direction_t d) {
     register uint64_t a7 asm("a7") = 177788;
     register uint64_t a0 asm("a0") = (uint64_t)d;
     asm volatile("ecall" : "+r"(a0) : "r"(a7) : "memory");
 }
\end{verbatim}
\end{codesnippet}

El minero (\texttt{miner.c}) contiene dos modos:

\begin{itemize}
  \item \emph{Cosecha básica}: avanzar, tomar cristales, volver a base (análogo a la tarea ejemplo del TP).
  \item \emph{Prueba de excepción}: cuando el minero tiene \texttt{id}$=$5, provoca una escritura en la VA \texttt{0x0}
  para generar un page fault (equivalente a la ``división por cero'' del TP3).
\end{itemize}
